\documentclass[12pt,a4paper]{article}

\usepackage[T1]{fontenc}
\usepackage[utf8]{inputenc}
\usepackage[brazil]{babel}
\usepackage{lmodern}
\usepackage{geometry}
\usepackage{hyperref}
\usepackage{enumitem}
\usepackage{array}

\geometry{margin=2.5cm}
\hypersetup{
  colorlinks=true,
  linkcolor=blue,
  urlcolor=blue
}

\title{Plano de Aula: Aula 3 -- Desenvolvimento Ágil}
\date{}

\begin{document}
\maketitle

\noindent
\textbf{Disciplina:} Engenharia de Software (Curso Técnico em Desenvolvimento de Sistemas)\\
\textbf{Duração:} 2 horas (120 minutos)\\
\textbf{Modalidade:} Presencial\\
\textbf{Materiais da aula:} \href{aula-03.html}{slides (HTML)} \quad \href{aula-03.pdf}{ementa (PDF)}\\
\textbf{Livro-base (referência):} \href{../99-Materiais/Engenharia%20de%20software%20projetos%20e%20processos%20by%20Wilson%20de%20P%C3%A1dua%20Paula%20Filho%20.epub.pdf}{Engenharia de software: projetos e processos (PDF)}

\section{Competências e Habilidades}
\textbf{Competências}
\begin{itemize}[leftmargin=*,itemsep=0.25em]
  \item Utilizar metodologias ágeis para gerenciar projetos de software.
  \item Aplicar princípios de engenharia de software para lidar com mudanças e reduzir riscos.
\end{itemize}

\textbf{Habilidades}
\begin{itemize}[leftmargin=*,itemsep=0.25em]
  \item Compreender desenvolvimento iterativo e incremental.
  \item Montar backlog e definir critérios de aceitação.
  \item Entender práticas básicas de Scrum e Kanban.
\end{itemize}

\section{Objetivos da Aula}
\begin{itemize}[leftmargin=*,itemsep=0.25em]
  \item Diferenciar iterativo vs.\ incremental e relacionar com métodos ágeis.
  \item Compreender conceitos centrais de Scrum e Kanban.
  \item Criar backlog simples e priorizar por valor e risco.
  \item Aplicar WIP limit e discutir efeitos da multitarefa.
\end{itemize}

\section{Assuntos Abordados no Dia}
\begin{itemize}[leftmargin=*,itemsep=0.25em]
  \item Agilidade como resposta a mudanças: ciclos curtos e feedback.
  \item Iterativo x incremental: definição e exemplos.
  \item Backlog: itens, priorização e critérios de aceitação.
  \item Scrum: eventos e objetivos (sprint, daily, review, retro).
  \item Kanban: fluxo, quadro e limites de trabalho em andamento (WIP).
\end{itemize}

\section{Metodologia}
Aula com exposição dialogada e demonstrações (quadros Scrum/Kanban). A atividade central é construir backlog e simular fluxo Kanban, reforçando definição de pronto e critérios de aceitação. O professor intervém com perguntas de clareza (o que é ``pronto''?) e de priorização (por que este item vem antes?).

\section{Materiais e Recursos}
\begin{itemize}[leftmargin=*,itemsep=0.25em]
  \item Quadro e marcadores (ou cartolina).
  \item Post-its (ideal) ou cartões de papel para montar backlog e quadro Kanban.
  \item Slides: \href{aula-03.html}{aula-03.html} e ementa: \href{aula-03.pdf}{aula-03.pdf}.
  \item Livro-base (PDF): \href{../99-Materiais/Engenharia%20de%20software%20projetos%20e%20processos%20by%20Wilson%20de%20P%C3%A1dua%20Paula%20Filho%20.epub.pdf}{Engenharia de software: projetos e processos}.
  \item Vídeos complementares (links nos slides).
\end{itemize}

\section{Cronograma e Descrição Detalhada (120 Minutos)}
\subsection*{Parte 1: Abertura e contexto (10 min)}
\begin{itemize}[leftmargin=*,itemsep=0.35em]
  \item Retomar processo/ciclo de vida (Aula 2) e motivar agilidade com mudança.
  \item Apresentar objetivos e dinâmica prática do encontro.
\end{itemize}

\subsection*{Parte 2: Iterativo x incremental (20 min)}
\begin{itemize}[leftmargin=*,itemsep=0.35em]
  \item Explicar diferença e dar exemplos do cotidiano de produto.
  \item Conectar com redução de risco: entregar cedo e aprender cedo.
\end{itemize}

\subsection*{Parte 3: Backlog e critérios de aceitação (25 min)}
\begin{itemize}[leftmargin=*,itemsep=0.35em]
  \item Demonstrar um backlog: itens, valor, prioridade, dependências.
  \item Critérios de aceitação: por que reduzem divergência e retrabalho.
  \item Exercício guiado: escrever 2 itens de backlog com aceitação.
\end{itemize}

\subsection*{Parte 4: Scrum e Kanban (25 min)}
\begin{itemize}[leftmargin=*,itemsep=0.35em]
  \item Scrum: sprint, daily, review e retro (objetivo de cada evento).
  \item Kanban: colunas, fluxo e WIP limit; diferença entre ``começar'' e ``terminar''.
  \item Mostrar um quadro e simular movimentação de itens.
\end{itemize}

\subsection*{Parte 5: Atividade prática em grupo (30 min)}
\begin{itemize}[leftmargin=*,itemsep=0.35em]
  \item Em grupos, criar 8 itens de backlog para o ``app organizador''.
  \item Priorizar 3 itens e definir critérios de aceitação.
  \item Montar quadro Kanban e mover 2 itens até concluído, respeitando WIP.
  \item Registrar lições: onde travou, por quê, o que mudaria no processo.
\end{itemize}

\subsection*{Parte 6: Fechamento (10 min)}
\begin{itemize}[leftmargin=*,itemsep=0.35em]
  \item Compartilhar conclusões e conectar com requisitos (Aula 4).
  \item Síntese: processo + método + ferramenta, com critérios e visibilidade.
\end{itemize}

\section{Avaliação}
\begin{itemize}[leftmargin=*,itemsep=0.25em]
  \item Participação na atividade e clareza dos itens de backlog.
  \item Qualidade dos critérios de aceitação e justificativa de priorização.
\end{itemize}

\section{Referências}
\begin{itemize}[leftmargin=*,itemsep=0.25em]
  \item \href{../99-Materiais/Engenharia%20de%20software%20projetos%20e%20processos%20by%20Wilson%20de%20P%C3%A1dua%20Paula%20Filho%20.epub.pdf}{PAULA FILHO, Wilson de Pádua. Engenharia de software: projetos e processos. 4. ed. 2019.}
\end{itemize}

\end{document}
